\title{Controlling Text-to-Image Diffusion by Orthogonal Finetuning}

\begin{abstract}
Large-scale text-to-image diffusion models exhibit remarkable ability in producing photorealistic images based on text inputs. Effectively steering or controlling these robust models to handle various downstream tasks remains a significant open challenge. To address this, we propose a systematic finetuning approach—Orthogonal Finetuning~(OFT)—for adapting text-to-image diffusion models to specific downstream applications. Unlike current techniques, OFT can theoretically maintain hyperspherical energy, which describes the pairwise relationships between neurons on a unit hypersphere. We observe that maintaining this property is essential for preserving the semantic generation capabilities of text-to-image diffusion models. To enhance finetuning stability, we introduce Constrained Orthogonal Finetuning (COFT), which enforces an additional radius constraint on the hypersphere. We focus on two key finetuning tasks for text-to-image models: subject-driven generation, where the objective is to create subject-specific images from a limited set of subject images and a text prompt, and controllable generation, where the model is designed to accept extra control signals. Our empirical results demonstrate that the OFT framework surpasses existing methods in terms of both generation quality and convergence speed.
\end{abstract}

\section{Introduction}

Recent advances in text-to-image diffusion models~\cite{saharia2022photorealistic,ramesh2022hierarchical,rombach2022high} have led to impressive capabilities in text-guided high-fidelity image generation. Despite these advances, text prompts alone can remain ambiguous and often fall short in providing precise and detailed control over the generated images. To mitigate this limitation, we concentrate on two categories of text-to-image generation tasks in this work:

\begin{itemize}[leftmargin=*,nosep]

    \item \textbf{Subject-driven generation}~\cite{ruiz2023dreambooth}: The goal is to produce images of a specific subject placed in new contexts, given only a handful of images of that subject along with a guiding text prompt.
    \item \textbf{Controllable generation}~\cite{zhang2023adding,mou2023t2i}: The aim is to generate images that adhere to additional control signals (such as canny edges or segmentation maps) in combination with a text prompt.
\end{itemize}

Both tasks fundamentally revolve around how to efficiently finetune text-to-image diffusion models without compromising their pretrained generative capabilities. We outline the essential criteria for an effective finetuning approach as: (1) \emph{training efficiency}: minimizing the number of trainable parameters and training epochs, and (2) \emph{generalizability preservation}: maintaining high-fidelity and diverse generative outputs. To achieve this, finetuning is generally performed either by adjusting neuron weights with a small learning rate (\emph{e.g.}, \cite{ruiz2023dreambooth}) or by incorporating a small module with re-parameterized neuron weights (\emph{e.g.}, \cite{hulora2022,zhang2023adding}). Despite their simplicity, neither approach can guarantee retention of the pretrained generative performance. Furthermore, there is a lack of a systematic framework for devising an effective finetuning strategy and selecting appropriate hyperparameters such as the number of training epochs. A major challenge lies in the absence of a metric to quantify how well the pretrained generative ability is preserved. Current finetuning methods implicitly assume that a smaller Euclidean distance between the finetuned model and the pretrained model corresponds to better preservation of pretrained capabilities. For this reason, finetuning typically employs very small learning rates. While this assumption may sometimes hold true, we contend that the Euclidean distance alone cannot fully reflect the extent of semantic preservation. Thus, a more structural metric to characterize the differences between the finetuned and pretrained models could significantly improve both the retention of pretrained performance and the stability of finetuning.

Motivated by empirical findings that hyperspherical similarity effectively encodes semantic information~\cite{liu2017deep,liu2018decoupled,chen2020angular}, we employ hyperspherical energy~\cite{liu2018learning} to describe the pairwise relational structure among neurons. Hyperspherical energy is defined as the aggregate hyperspherical similarity (\emph{e.g.}, cosine similarity) between all pairs of neurons within the same layer, reflecting the degree of neuron uniformity on the unit hypersphere~\cite{liu2021learning}. We propose that a well-finetuned model should exhibit minimal deviation in hyperspherical energy relative to the pretrained model. A straightforward approach might involve adding a regularizer to maintain constant hyperspherical energy during finetuning, but this does not guarantee effective minimization of the energy difference. To address this, we leverage an invariance property of hyperspherical energy — the pairwise hyperspherical similarity is theoretically preserved when the same orthogonal transformation is applied to all neurons. Inspired by this invariance, we introduce Orthogonal Finetuning~(OFT), which adapts large text-to-image diffusion models to downstream tasks without altering their hyperspherical energy. The core concept is to learn a layer-shared orthogonal transformation for neurons that maintains their pairwise angles. OFT can also be interpreted as redefining the canonical coordinate system for neurons within the same layer. Considering that a smaller Euclidean distance between the finetuned and pretrained models corresponds to better retention of pretrained performance, we additionally propose a variant called Constrained Orthogonal Finetuning~(COFT), which restricts the finetuned model to lie within a hypersphere of fixed radius centered around the pretrained neurons.

The reasoning behind why orthogonal transformation is effective for finetuning neurons partially originates from the 2D Fourier transform, which decomposes an image into its magnitude and phase spectra. The phase spectrum, representing the angular relationship between the input and the basis, retains the majority of semantic information. For instance, an image’s phase spectrum combined with a random magnitude spectrum can still reconstruct the original image without losing its semantic content. This observation implies that altering the directions of neurons is crucial for semantically modifying the generated image, which aligns with the objectives of both subject-driven and controllable generation. However, allowing unrestricted changes to neuron directions would inevitably compromise the generative performance established during pretraining. To limit this freedom, we propose maintaining the angles between every pair of neurons, based on the hypothesis that these inter-neuronal angles are fundamental to encoding the neural network’s knowledge. Guided by this intuition, it is natural to learn a layer-shared orthogonal transformation for neurons in each layer to keep the hyperspherical energy constant.

Our approach is also inspired by orthogonal over-parameterized training~\cite{liu2021orthogonal}, which trains classification neural networks from scratch by applying orthogonal transformations to a randomly initialized network. This is motivated by the fact that a randomly initialized network is expected to have low hyperspherical energy, and \cite{liu2021orthogonal} aims to maintain this low energy during training (as lower energy correlates with better classification generalization~\cite{liu2018learning,lin2020regularizing}). The work in \cite{liu2021orthogonal} demonstrates that orthogonal transformations provide sufficient flexibility to train neural networks with good generalization for classification tasks. In contrast, our focus is on finetuning text-to-image diffusion models to enhance controllability and improve downstream generative capabilities. We highlight two main distinctions between OFT and \cite{liu2021orthogonal}: First, whereas \cite{liu2021orthogonal} minimizes hyperspherical energy, OFT is designed to preserve the pretrained model’s hyperspherical energy, ensuring that the original pretrained structure remains intact during finetuning. Minimizing hyperspherical energy in diffusion model finetuning could disrupt the established semantic structures. Second, OFT aims to minimize deviation from the pretrained model, resulting in a constrained variant, whereas \cite{liu2021orthogonal} does not impose such constraints. Achieving a balance between flexibility and stability is essential for effective finetuning, and we contend that our OFT framework successfully accomplishes this. Our contributions are summarized as follows:

\begin{itemize}[leftmargin=*,nosep]

    \item We introduce a novel finetuning approach called Orthogonal Finetuning (OFT) to enhance the controllability of text-to-image diffusion models. To further boost stability, we also propose a constrained variant that restricts the angular deviation from the pretrained model.
    \item Unlike existing finetuning methods, OFT adjusts the model parameters while theoretically ensuring the preservation of hyperspherical energy, which we empirically identify as a crucial factor for maintaining the generative semantic integrity of the pretrained model.
    \item We evaluate OFT on two applications: subject-driven generation and controllable generation. Through extensive empirical experiments, we show substantial improvements over prior approaches in terms of generation quality, convergence speed, and finetuning stability. Furthermore, OFT demonstrates higher sample efficiency, achieving good convergence with significantly fewer training images and epochs.
    \item For controllable generation, we propose a novel control consistency metric to measure controllability. The key idea is to infer the control signal from the generated image and then compare it to the original control signal. This metric further confirms the strong controllability achieved by OFT.
\end{itemize}

\section{Related Work}

\textbf{Text-to-image diffusion models}. Significant advancements~\cite{saharia2022photorealistic,ramesh2022hierarchical,rombach2022high,gu2022vector,nichol2022glide} have been achieved in text-to-image synthesis, largely due to the rapid progress in diffusion-based generative models~\cite{ho2020denoising,song2021score,song2021denoising,dhariwal2021diffusion} and vision-language representation learning~\cite{su2020vlbert,lu2019vilbert,radford2021learning,wang2021simvlm,li2022blip,li2023blip,alayrac2022flamingo,schuhmann2022laion}. GLIDE~\cite{nichol2022glide} and Imagen~\cite{saharia2022photorealistic} train diffusion models directly in pixel space. GLIDE jointly trains the text encoder with a diffusion prior using paired text-image data, whereas Imagen relies on a frozen pretrained text encoder. Stable Diffusion~\cite{rombach2022high} and DALL-E2~\cite{ramesh2022hierarchical} operate in latent spaces; Stable Diffusion employs VQ-GAN~\cite{esser2021taming} to learn a visual codebook as the latent space, while DALL-E2 uses CLIP~\cite{radford2021learning} to construct a joint embedding space for images and text. Beyond diffusion models, generative adversarial networks~\cite{reed2016generative,zhang2017stackgan,xu2018attngan,li2019controllable} and autoregressive models~\cite{ramesh2021zero,ding2021cogview,wu2022nuwa,yuscaling} have also been explored for text-to-image generation. OFT is inherently model-agnostic and can be applied to any text-to-image diffusion model.

\textbf{Subject-driven generation}. To avoid altering the subject, \cite{avrahami2022blended,nichol2022glide} incorporate user-provided masks as additional conditions. Inversion techniques~\cite{choi2021ilvr,dhariwal2021diffusion,ramesh2022hierarchical,gal2022image} allow for context modification without changing the subject. \cite{hertz2022prompt} supports both local and global edits without needing input masks. However, these approaches struggle to maintain identity-specific details of the subject. Pivotal Tuning~\cite{roich2022pivotal} finetunes the generator around an initially inverted latent code with regularization to retain identity. Similarly, \cite{nitzan2022mystyle} learns a personalized generative face prior from multiple images of an individual. \cite{casanova2021instance} produces different variations of an instance but may sacrifice instance-specific details. DreamBooth~\cite{ruiz2023dreambooth} finetunes a text-to-image diffusion model with a small set of subject images and a customized token, using reconstruction and class-specific prior preservation losses. OFT builds upon the DreamBooth framework but replaces naive finetuning with orthogonal transformations.

\textbf{Controllable generation}. Image-to-image translation can be regarded as a type of controllable generation, with many earlier approaches utilizing conditional generative adversarial networks~\cite{isola2017image,zhu2017toward,wang2018high,park2019semantic,choi2018stargan}. Diffusion models have also been employed for image-to-image translation tasks~\cite{saharia2022palette,voynov2022sketch,wang2022pretraining}. Recently, ControlNet~\cite{zhang2023adding} introduced a method to steer a pretrained diffusion model by fine-tuning and adapting it to incorporate additional control signals, achieving impressive controllable generation results. Another similar and contemporary approach, T2I-Adapter~\cite{mou2023t2i}, likewise fine-tunes pretrained diffusion models to enhance controllability over the generated images. Building on the same task framework as in \cite{zhang2023adding,mou2023t2i}, we employ OFT to fine-tune pretrained diffusion models, resulting in consistently improved controllability while requiring less training data and fewer finetuning parameters. Importantly, OFT adds no extra computational cost during inference.

\textbf{Model finetuning}. The practice of fine-tuning large pretrained models for downstream applications has become increasingly widespread~\cite{devlin2018bert,brown2020language,he2022masked}. Within this context, adaptation techniques (\emph{e.g.}, \cite{hulora2022,houlsby2019parameter,pfeiffer2020adapterfusion}) have been extensively explored in natural language processing. LoRA~\cite{hulora2022} is the most closely related method to OFT, as it models the finetuning weight update as a low-rank additive modification. In contrast, OFT employs a layer-shared orthogonal transformation to adjust neuron weights through a multiplicative update, which provably maintains the pairwise angles between neurons in the same layer, enhancing stability.

\section{Orthogonal Finetuning}

\subsection{Why Does Orthogonal Transformation Make Sense?}

We begin by examining why applying orthogonal transformation is beneficial for finetuning text-to-image diffusion models. This inquiry can be divided into two key questions: (1) why adjusting the angles (i.e., directions) of neurons is desirable during finetuning, and (2) why orthogonal transformations are used to refine these angles.

Regarding the first question, we draw motivation from the empirical findings in \cite{liu2018decoupled,chen2020angular} that angular differences in features effectively represent the semantic gap. SphereNet~\cite{liu2017deep} demonstrates that training a neural network with all neurons normalized on a unit hypersphere provides comparable capacity and even enhances generalizability, suggesting that neuron directions can fully encapsulate the most critical information from data. To further highlight the significance of neuron angles, we conduct a simple experiment shown in Figure~\ref{fig:toy_ae}, where a standard convolutional autoencoder is trained on a set of flower images. During training, the standard inner product is used to produce the feature map ($z$ represents the output element of the convolution kernel $\bm{w}$, and $\bm{x}$ is the input in the sliding window). In testing, we compare three methods for generating the feature map: (a) the inner product used during training, (b) the magnitude information, and (c) the angular information. Results in Figure~\ref{fig:toy_ae} indicate that angular information of neurons can nearly perfectly reconstruct the input images, whereas neuron magnitudes provide no useful information. We stress that cosine activation (c) was not applied during training, which was exclusively based on the inner product. These results suggest that neuron angles (directions) predominantly store the semantic content of the input images. Therefore, finetuning neuron directions is likely to be more effective for altering semantic information.

For the second question, the most straightforward approach to finetune neuron directions is to simultaneously rotate or reflect all neurons within the same layer, naturally leading to orthogonal transformations. While other angular transformations allowing different neurons to rotate by varying angles might offer more flexibility, orthogonal transformation strikes a balance between flexibility and structural regularity. Additionally, \cite{liu2021orthogonal} demonstrates that orthogonal transformations possess sufficient expressive power for neural network learning. To support this claim, we conduct an experiment illustrating the effective regularization induced by the orthogonality constraint. Using the controllable generation setup from ControlNet~\cite{zhang2023adding}, results presented in Figure~\ref{fig:ortho} show that our standard OFT performs robustly and achieves precise control after training completes (epoch 20). In contrast, OFT without the orthogonality constraint fails to produce any realistic images and lacks any control capability. This experiment confirms the critical role of the orthogonality constraint in OFT.

\subsection{General Framework}

The standard finetuning approach generally employs gradient descent with a small learning rate to update a model (or specific layers within it). This small learning rate naturally promotes minimal changes from the pretrained model, and conventional finetuning essentially focuses on training the model while implicitly minimizing $\thickmuskip=2mu \medmuskip=2mu \| \bm{M} - \bm{M}^0\|$, where $\bm{M}$ represents the finetuned model parameters and $\bm{M}^0$ denotes the pretrained model parameters. Although this implicit restriction is intuitive, it may still allow too much flexibility when finetuning large models. To tackle this issue, LoRA introduces an extra low-rank constraint on the weight updates, namely, $\thickmuskip=2mu \medmuskip=2mu \text{rank}(\bm{M} - \bm{M}^0) = r'$, where $r'$ is chosen to be a small integer. Unlike LoRA, OFT imposes a constraint on the pairwise similarity among neurons: $\thickmuskip=2mu \medmuskip=2mu \| \text{HE}(\bm{M}) - \text{HE}(\bm{M}^0) \| = 0$, where $\text{HE}(\cdot)$ stands for the hyperspherical energy of a weight matrix. To illustrate, consider a fully connected layer $\thickmuskip=2mu \medmuskip=2mu \bm{W} = \{\bm{w}_1, \ldots, \bm{w}_n\} \in \mathbb{R}^{d \times n}$, where each neuron $\bm{w}_i \in \mathbb{R}^d$ (and $\bm{W}^0$ denotes the pretrained weights). The output vector $\thickmuskip=2mu \medmuskip=2mu \bm{z} \in \mathbb{R}^n$ from this layer is computed by $\thickmuskip=2mu \medmuskip=2mu \bm{z} = \bm{W}^\top \bm{x}$, with $\thickmuskip=2mu \medmuskip=2mu \bm{x} \in \mathbb{R}^d$ as the input vector. OFT can be viewed as minimizing the difference in hyperspherical energy between the finetuned and pretrained models:

\begin{equation}\label{eq:oft_principle}
\footnotesize
\min_{\bm{W}} \left\| \text{HE}(\bm{W}) - \text{HE}(\bm{W}^0) \right\|~~\Leftrightarrow~~ \min_{\bm{W}} \bigg{\|} \sum_{i \neq j} \| \hat{\bm{w}}_i - \hat{\bm{w}}_j \|^{-1} - \sum_{i \neq j} \| \hat{\bm{w}}_i^0 - \hat{\bm{w}}_j^0 \|^{-1} \bigg{\|}
\end{equation}

Here, $\thickmuskip=2mu \medmuskip=2mu \hat{\bm{w}}_i := \bm{w}_i / \|\bm{w}_i\|$ denotes the $i$-th normalized neuron, and the hyperspherical energy of the layer $\bm{W}$ is defined as $\thickmuskip=2mu \medmuskip=2mu \text{HE}(\bm{W}) := \sum_{i \neq j} \|\hat{\bm{w}}_i - \hat{\bm{w}}_j\|^{-1}$. It is straightforward to see that the minimum achievable value for Eq.~\eqref{eq:oft_principle} is zero. This minimum is attained when $\bm{W}$ and $\bm{W}^0$ differ only by a rotation or reflection, i.e., $\thickmuskip=2mu \medmuskip=2mu \bm{W} = \bm{R} \bm{W}^0$, where $\thickmuskip=2mu \medmuskip=2mu \bm{R} \in \mathbb{R}^{d \times d}$ is an orthogonal matrix (with determinant $1$ or $-1$ corresponding to rotation or reflection, respectively). This principle underlies OFT, which asserts that finetuning can

\begin{equation}\label{eq:oft_general}
    \footnotesize
    \bm{z}=\bm{W}^\top\bm{x}=(\bm{R}\cdot\bm{W}^0)^\top\bm{x},~~~\text{s.t.}~\bm{R}^\top\bm{R}=\bm{R}\bm{R}^\top=\bm{I}
\end{equation}

Here, $\bm{W}$ represents the OFT-finetuned weight matrix and $\bm{I}$ stands for the identity matrix. The OFT framework is depicted in Figure~\ref{fig:block}. Analogous to zero initialization in LoRA, it is essential to guarantee that OFT fine-tunes the pretrained model precisely starting from $\bm{W}^0$. To accomplish this, the orthogonal matrix $\bm{R}$ is initialized as the identity matrix, ensuring the finetuned model begins with the pretrained weights. To maintain the orthogonality of $\bm{R}$, differential orthogonalization methods discussed in \cite{liu2021orthogonal,lezcano2019cheap} can be applied. We will further explain how to enforce orthogonality efficiently in terms of computation.

\subsection{Efficient Orthogonal Parameterization}\label{sect:eff_ortho}

Although standard orthogonalization techniques such as the Gram-Schmidt process are differentiable, they tend to be computationally intensive in practice~\cite{liu2021orthogonal}. To improve efficiency, we utilize the Cayley parameterization to construct the orthogonal matrix. Specifically, the orthogonal matrix is formed as $\thickmuskip=2mu \medmuskip=2mu \bm{R}=(\bm{I}+\bm{Q})(\bm{I}-\bm{Q})^{-1}$, where $\bm{Q}$ is a skew-symmetric matrix satisfying $\thickmuskip=2mu \medmuskip=2mu \bm{Q}=-\bm{Q}^\top$. The computational gains come with a minor limitation—the Cayley parameterization generates only orthogonal matrices with determinant $1$, belonging to the special orthogonal group. Fortunately, this constraint does not impair practical performance. Despite employing the Cayley transform for parameterizing the orthogonal matrix, $\bm{R}$ can still be parameter-inefficient for large $d$. To mitigate this, we propose representing $\bm{R}$ as a block-diagonal matrix comprising $r$ blocks, resulting in the following structure:

\begin{equation}
    \footnotesize
    \bm{R}=\text{diag}(\bm{R}_1,\bm{R}_2,\cdots,\bm{R}_r)=\begin{bmatrix}
    \bm{R}_1\in \text{O}(\frac{d}{r}) &  &\\
    &\ddots & \\
    & & \bm{R}_r\in \text{O}(\frac{d}{r})
    \end{bmatrix}\in \text{O}(d)
\end{equation}

where $\text{O}(d)$ represents the orthogonal group in dimension $d$, and $\thickmuskip=2mu \medmuskip=2mu \bm{R}\in\mathbb{R}^{d\times d}$ and $\thickmuskip=2mu \medmuskip=2mu \bm{R}_i\in\mathbb{R}^{d/r\times d/r},\forall i$ are orthogonal matrices. When $\thickmuskip=2mu \medmuskip=2mu r=1$, the block-diagonal orthogonal matrix reduces to a conventional unconstrained orthogonal matrix. For an orthogonal matrix sized $\thickmuskip=2mu \medmuskip=2mu d\times d$, there are $\thickmuskip=2mu \medmuskip=2mu d(d-1)/2$ parameters, leading to a complexity of $\mathcal{O}(d^2)$. In the case of an $r$-block diagonal orthogonal matrix, the parameter count is $\thickmuskip=2mu \medmuskip=2mu d(d/r-1)/2$, resulting in complexity $\mathcal{O}(d^2/r)$. Optionally, we may enforce block sharing such that $\thickmuskip=2mu \medmuskip=2mu \bm{R}_i=\bm{R}_j,\forall i\neq j$, which further diminishes parameter count to $\mathcal{O}(d^2/r^2)$. Despite these parameter efficiency strategies, the resulting matrix $\bm{R}$ remains orthogonal, thus fully preserving hyperspherical energy.

We compare the parameter efficiency of OFT against LoRA. For LoRA with low-rank parameter $r'$, the number of trainable parameters is $\thickmuskip=2mu \medmuskip=2mu r'(d+n)$. Assuming both $r$ and $r'$ scale with the neuron dimension $d$ (e.g., $\thickmuskip=2mu \medmuskip=2mu r=r'=\alpha d$, with some constant $0<\alpha\leq 1$), LoRA's parameter complexity becomes $\mathcal{O}(d^2+dn)$ while OFT achieves parameter complexity $\mathcal{O}(d)$. To concretely illustrate this difference, consider a weight matrix sized $\thickmuskip=2mu \medmuskip=2mu 128\times 128$: LoRA with $r'=8$ has $2,048$ trainable parameters, whereas OFT with $r=8$ (no block sharing) has only $960$ trainable parameters.

\subsection{Constrained Orthogonal Finetuning}

We can restrict the flexibility of the original OFT by enforcing the finetuned model to remain close to the pretrained model within a small neighborhood. Specifically, COFT employs the forward pass:

\begin{equation}\label{eq:coft_general}
    \footnotesize
    \bm{z}=\bm{W}^\top\bm{x}=(\bm{R}\cdot\bm{W}^0)^\top\bm{x},~~~\text{s.t.}~\bm{R}^\top\bm{R}=\bm{R}\bm{R}^\top=\bm{I},~\norm{\bm{R}-\bm{I}}\leq \epsilon
\end{equation}

which imposes both an orthogonality constraint and a constraint limiting $\bm{R}$’s deviation from the identity matrix by $\epsilon$. The orthogonality condition can be enforced via the Cayley parameterization introduced in Section~\ref{sect:eff_ortho}. However, incorporating the $\epsilon$-deviation constraint into the Cayley-parameterized orthogonal matrix is nontrivial. To better understand the Cayley transform, we approximate $\thickmuskip=2mu \medmuskip=2mu \bm{R}=(\bm{I}+\bm{Q})(\bm{I}-\bm{Q})^{-1}$ using the Neumann series as $\thickmuskip=2mu \medmuskip=2mu \bm{R}\approx\bm{I}+2\bm{Q}+\mathcal{O}(\bm{Q}^2)$ (assuming the Neumann series converges).

(series converges in the operator norm). Consequently, we can incorporate the constraint $\thickmuskip=2mu \medmuskip=2mu\|\bm{R}-\bm{I}\|\leq\epsilon$ within the Cayley transform, making the equivalent constraint $\thickmuskip=2mu \medmuskip=2mu \|\bm{Q}-\bm{0}\|\leq\epsilon'$, where $\bm{0}$ represents an all-zero matrix and $\epsilon'$ is a different error hyperparameter (distinct from $\epsilon$). This new constraint on matrix $\bm{Q}$ can be easily enforced using projected gradient descent. To initialize the orthogonal matrix $\bm{R}$ as the identity, we set $\bm{Q}$ to an all-zero matrix initially. COFT can be interpreted as combining two explicit constraints: minimizing hyperspherical energy difference and restricting deviation from the pretrained model. The latter constraint is commonly implicit in existing finetuning methods, but COFT explicitly imposes it. Although OFT shows excellent performance, we observe that COFT provides even greater finetuning stability due to this explicit deviation constraint. Figure~\ref{fig:eps_coft} illustrates how $\epsilon$ influences COFT's performance. It shows that $\epsilon$ regulates finetuning flexibility: larger $\epsilon$ makes the COFT-finetuned model resemble the OFT-finetuned model, whereas smaller $\epsilon$ causes it to behave more like the pretrained text-to-image diffusion model.

\subsection{Re-scaled Orthogonal Finetuning}

We propose a straightforward extension of the original OFT by additionally learning a magnitude scaling factor for each neuron. This is motivated by the observation that re-scaling neurons does not affect the hyperspherical energy since magnitude normalization removes its influence. Specifically, the forward pass is given by: $\thickmuskip=2mu \medmuskip=2mu\bm{z}=(\bm{R}\bm{W}^0\bm{D})^\top\bm{x}$\footnote{\textbf{Errata}: In the NeurIPS camera-ready version, the forward pass of re-scaled OFT was incorrectly stated as $\thickmuskip=2mu \medmuskip=2mu\bm{z}=(\bm{D}\bm{R}\bm{W}^0)^\top\bm{x}$. The original implementation was correct, so the results in Appendix~\ref{app:rescaled} remain valid.} where $\thickmuskip=2mu \medmuskip=2mu \bm{D}=\text{diag}(s_1,\cdots,s_n)\in\mathbb{R}^{n\times n}$ is a learnable diagonal matrix with positive diagonal entries $s_1,\cdots,s_n$. Unlike the original OFT forward pass in Eq.~\eqref{eq:oft_general}, where only $\bm{R}$ is trainable, here both the diagonal matrix $\bm{D}$ and the orthogonal matrix $\bm{R}$ are learned. This re-scaled OFT enhances the flexibility of OFT while introducing a negligible number of extra parameters. Although we use the original OFT in experiments to demonstrate the effectiveness of orthogonal transformation alone, we find that re-scaled OFT generally performs better (see Appendix~\ref{app:rescaled}).

\section{Intriguing Insights and Discussions}

\textbf{OFT is architecture-agnostic}. In principle, OFT can be applied to any neural network architecture. For Transformers, LoRA is usually applied to attention weights~\cite{hulora2022}. To ensure a fair comparison with LoRA, we restrict OFT finetuning to attention weights in our experiments. Beyond fully connected layers, OFT is also well-suited for finetuning convolutional layers due to the block-diagonal structure of $\bm{R}$ having meaningful interpretations in convolutional contexts (unlike LoRA). When the number of blocks matches the number of input channels, each block transforms a distinct neuron channel, analogous to learning depth-wise convolution kernels~\cite{chollet2017xception}. When all blocks in $\bm{R}$ are shared, OFT applies the same orthogonal transformation across channels. We conduct a preliminary investigation of OFT finetuning convolutional layers in Appendix~\ref{app:convolution}.

\textbf{Relation to LoRA}. LoRA adds a low-rank matrix to ensure the pretrained weight matrix's information does not change drastically. In contrast, OFT enforces an orthogonal (full-rank) transformation on the pretrained weight matrix, thereby preserving the pretrained knowledge. The forward pass of OFT can be expressed as $\thickmuskip=2mu \medmuskip=2mu \bm{z}=(\bm{R}\bm{W}^0)^\top\bm{x}=(\bm{W}^0+(\bm{R}-\bm{I})\bm{W}^0)^\top\bm{x}$, where $\thickmuskip=2mu \medmuskip=2mu (\bm{R}-\bm{I})\bm{W}^0$ resembles the low-rank weight update in LoRA. Given that $\bm{W}^0$ is usually full-rank, OFT effectively implements a low-rank weight update when $\thickmuskip=2mu \medmuskip=2mu \bm{R}-\bm{I}$ has low rank. Similarly to LoRA’s rank hyperparameter $r'$, OFT uses a diagonal block size parameter $r$ to limit the number of trainable parameters. Notably, LoRA and OFT adopt fundamentally different approaches to parameter efficiency: LoRA leverages low-rank structure for parameter reduction, whereas OFT exploits sparsity via block-diagonal orthogonality.

\textbf{Reasons for faster OFT convergence}. As illustrated in Figure~\ref{fig:toy_ae}, the most effective semantic update corresponds to altering neuron directions, which aligns with OFT’s design goal. Moreover, OFT can be interpreted as finetuning neurons on a smooth hyperspherical manifold, leading to improved optimization landscapes. This advantage is supported empirically by \cite{liu2021orthogonal}.

\textbf{Why hyperspherical energy is not minimized}. Unlike \cite{liu2021orthogonal}, our approach does not seek to minimize hyperspherical energy. In classification tasks, non-redundant neurons are preferred. Minimizing hyperspherical energy results in neurons being uniformly distributed on the hypersphere, which is not a suitable objective for finetuning, as it could disrupt the pretrained representations.

\textbf{Balance between flexibility and regularization in finetuning}. We identify a fundamental trade-off between flexibility and regularity. Standard finetuning offers the greatest flexibility but often results in poor stability and can easily lead to model collapse. Surprisingly simple, OFT achieves a good compromise by maintaining the pairwise neuron angles. Additionally, the block-diagonal parameterization can be interpreted as a stronger form of regularization on the orthogonal matrix.

\textbf{No extra inference cost}. Unlike ControlNet, our OFT framework does not add any additional inference overhead to the finetuned model. During inference, we simply multiply the learned orthogonal matrix $\bm{R}$ with the pretrained weight matrix $\bm{W}^0$ to produce an equivalent weight matrix $\thickmuskip=2mu \medmuskip=2mu \bm{W}=\bm{R}\bm{W}^0$. Therefore, the inference speed remains the same as that of the pretrained model.

\section{Experiments and Results}

\textbf{Experimental setup}. In our experiments, we utilize Stable Diffusion v1.5~\cite{rombach2022high} as the pretrained text-to-image model. For fairness, we randomly select generated images from each method. For subject-driven generation, we generally follow the protocol of DreamBooth~\cite{ruiz2023dreambooth}. For controllable generation, we mostly adhere to the settings of ControlNet~\cite{zhang2023adding} and T2I-Adapter~\cite{mou2023t2i}. To ensure a fair comparison with LoRA, we apply OFT or COFT only to the same layer where LoRA is employed. Additional results and details are provided in Appendix~\ref{app:settings}.

\subsection{Subject-driven Generation}

\textbf{Setup}. We use DreamBooth~\cite{ruiz2023dreambooth} and LoRA~\cite{hulora2022} as baseline methods. All approaches employ the same loss function as DreamBooth. For DreamBooth and LoRA, we generally follow the original works and adopt their best hyperparameter configurations. Further results are available in Appendices~\ref{app:settings}, \ref{app:coft_oft}, \ref{app:more_qualitative}, and \ref{app:failure}.

\textbf{Finetuning stability and convergence}. We begin by assessing the finetuning stability and convergence rate for DreamBooth, LoRA, OFT, and COFT. The outcomes are presented in Figure~\ref{fig:teaser} and Figure~\ref{fig:db_convergence}. It is evident that both COFT and OFT can finetune the diffusion model in a stable manner. After 400 iterations, DreamBooth and OFT variants both demonstrate effective control, whereas LoRA fails to maintain the subject’s identity. At 2000 iterations, DreamBooth starts producing collapsed images, and LoRA is unable to generate a yellow shirt, instead producing yellow fur. In contrast, OFT and COFT continue to provide stable and consistent control over the generated images. These findings confirm the rapid yet stable convergence offered by our OFT framework in subject-driven generation. We highlight that the robustness to the number of finetuning iterations is particularly valuable, as it reduces the need to carefully tune the iteration count for different subjects. Both OFT and COFT allow the use of relatively large iteration numbers without extensive tuning. Specifically, for COFT with a suitable $\epsilon$, setting both the learning rate and iteration number becomes straightforward.

\textbf{Quantitative comparison}. Following \cite{ruiz2023dreambooth}, we perform a quantitative evaluation to measure subject fidelity (DINO~\cite{caron2021emerging}, CLIP-I~\cite{radford2021learning}), text prompt fidelity (CLIP-T~\cite{radford2021learning}), and sample diversity (LPIPS~\cite{zhang2018unreasonable}). CLIP-I calculates the average pairwise cosine similarity of CLIP embeddings between generated and real images. DINO operates similarly to CLIP-I but utilizes ViT S/16 DINO embeddings. CLIP-T computes the average cosine similarity of CLIP embeddings between the text prompt and generated images. Additionally, we assess the average LPIPS cosine similarity among generated images of the same subject conditioned on the same text prompt. Table~\ref{table:dreambooth_final} demonstrates that COFT and OFT significantly outperform DreamBooth and LoRA on the DINO and CLIP-I metrics, while showing slightly improved or comparable results in prompt fidelity and diversity metrics. To reduce randomness, we finetune the same pretrained model 30 times with different random seeds for each method. These results indicate that our OFT framework not only delivers enhanced convergence and stability but also consistently achieves superior final performance.

\textbf{Qualitative comparison}. To better illustrate the advantages of OFT, we present several randomly selected examples of subject-driven generation in Figure~\ref{fig:dreambooth_final}. For an equitable comparison, all examples are produced using the same finetuned model for each method, ensuring that no text prompt is separately optimized for the final output. For every approach, we pick the model that attains the highest validation CLIP scores. The results shown in Figure~\ref{fig:dreambooth_final} indicate that both OFT and COFT maintain excellent semantic consistency of the subject, whereas LoRA frequently struggles to retain the subject's identity (for instance, LoRA entirely loses the subject identity in the bowl example). Additionally, OFT and COFT offer significantly more precise control through text prompts, while DreamBooth, despite preserving the subject identity, often fails to generate images that adhere to the text prompt (e.g., the first row in the bowl example). This qualitative comparison confirms that our OFT framework provides superior controllability alongside robust subject preservation. Furthermore, OFT's performance is stable across different numbers of iterations, delivering strong results even with a high iteration count, a property not shared by DreamBooth or LoRA. Additional qualitative examples are available in Appendix~\ref{app:more_qualitative}. We also performed a human evaluation detailed in Appendix~\ref{app:human}, which further supports the advantages of OFT.

\subsection{Controllable Generation}

\textbf{Settings}. We compare against ControlNet~\cite{zhang2023adding}, T2I-Adapter~\cite{mou2023t2i}, and LoRA~\cite{hulora2022} as baseline methods. This study focuses on three challenging controllable generation tasks: Canny edge to image (C2I) on the COCO dataset~\cite{lin2014microsoft}, segmentation map to image (S2I) on the ADE20K dataset~\cite{zhou2017scene}, and landmark to face (L2F) on the CelebA-HQ dataset~\cite{karras2018progressive,xia2021tedigan}. Each method is used to finetune Stable Diffusion (SD) v1.5 on these datasets for 20 epochs. Additional results can be found in Appendices~\ref{app:more_qualitative}, \ref{app:more_control_task}, and \ref{app:failure}.

\textbf{Convergence}. We assess the convergence speed of ControlNet, T2I-Adapter, LoRA, and COFT on the L2F task through both quantitative and qualitative analyses. For the evaluation metric, we calculate the average $\ell_2$ distance between the control face landmarks and the predicted face landmarks. In Figure~\ref{fig:face_convergence}, we display the face landmark error for models finetuned over different numbers of epochs. It is evident that both COFT and OFT achieve considerably faster convergence. LoRA requires 20 epochs to reach the performance level that our OFT framework attains by the 8th epoch. We also observe that OFT and COFT employ a comparable number of trainable parameters as LoRA (significantly fewer than ControlNet), yet they converge much more efficiently than existing approaches. Furthermore, the rapid convergence of OFT is supported by the findings in Figure~\ref{fig:teaser}. The right example in this figure demonstrates that OFT is notably more data-efficient compared to ControlNet and LoRA, as it can achieve strong convergence using only 5\% of the ADE20K dataset. For qualitative comparisons, we concentrate on OFT, COFT, and ControlNet, since ControlNet attains the closest landmark error to our methods. The outcomes shown in Figure~\ref{fig:face_convergence_qual} indicate that both OFT and COFT converge steadily, with the generated face pose progressively aligning with the control landmarks. In contrast, ControlNet exhibits a sudden onset of controllability after the 8th epoch, consistent with the abrupt convergence phenomenon reported in \cite{zhang2023adding}. This stability in convergence renders our OFT framework easier to apply in practice, as its training dynamics are significantly smoother and more predictable. Consequently, identifying optimal hyperparameters for OFT is simpler.

\textbf{Quantitative comparison}. We propose a control consistency metric to assess the effectiveness of controllable generation. The core concept is to extract the control signal from the generated image and then compare it against the original input control signal. For the C2I task, we calculate IoU and F1 scores. In the case of the S2I task, we measure mean IoU, mean accuracy, and overall accuracy. For the L2F task, we determine the mean $\ell_2$ distance between the control landmarks and the predicted landmarks. Additional details about these consistency metrics are provided in Appendix~\ref{app:settings}. For all methods evaluated, we apply the optimal hyperparameter configurations. The results presented in Table~\ref{table:control_gen} demonstrate that both OFT and COFT achieve significantly stronger and more precise control compared to other approaches. We note that adapter-based methods (such as T2I-Adapter and ControlNet) show slower convergence and inferior final outcomes. Relative to ControlNet, LoRA performs better on the S2I task but underperforms on the C2I and L2F tasks. Overall, we observe that LoRA’s performance ceiling is comparatively low, even with meticulous hyperparameter tuning. In contrast, the performance of our OFT framework appears not to have plateaued, as we empirically find it continues to improve with an increased number of trainable parameters. We highlight that our quantitative evaluation framework for controllable generation represents a novel contribution, as it provides an accurate assessment of the control capabilities of finetuned models (subject to the accuracy limits of the utilized off-the-shelf segmentation/detection model).

\textbf{Qualitative comparison}. We further conduct a qualitative comparison of OFT and COFT against leading state-of-the-art approaches such as ControlNet, T2I-Adapter, and LoRA. The randomly generated images shown in Figure~\ref{fig:control_final} illustrate that OFT and COFT not only produce high-fidelity and realistic images but also achieve precise control. In the S2I task, it is evident that LoRA fails entirely to generate images consistent with the input segmentation map, whereas ControlNet, OFT, and COFT effectively guide the image generation. Compared to ControlNet, both OFT and COFT generate higher fidelity images with richer details and more accurate geometric structures while utilizing significantly fewer model parameters. For the C2I task, OFT and COFT can hallucinate semantically coherent images from coarse Canny edges, whereas T2I-Adapter and LoRA perform notably worse. In the L2F task, our approach delivers the most precise pose control of generated faces, even under difficult facial poses. Across all three control tasks, OFT and COFT produce qualitatively superior images relative to the current best baselines, underscoring the efficacy of the OFT framework in controllable generation. To offer a more thorough qualitative comparison, additional examples for all three control tasks are provided in Appendix~\ref{app:control_exp}, and we also showcase the strong performance of OFT on further control tasks (such as dense pose to human body, sketch to image, and depth to image) in Appendix~\ref{app:more_control_task}.

\section{Concluding Remarks and Open Problems}

Inspired by the insight that angular relationships among neurons play a key role in defining visual semantics, we introduce a straightforward yet powerful finetuning technique—orthogonal finetuning (OFT)—for controlling text-to-image diffusion models. Our focus is on two applications: subject-driven generation and controllable generation. In comparison with prior approaches, OFT offers enhanced controllability and greater finetuning stability while using fewer finetuning parameters. Crucially, OFT does not add any extra computational cost during inference, making the resulting model highly efficient for deployment.

OFT also raises several intriguing open questions. First, the orthogonality constraint in OFT is enforced through Cayley parametrization, which requires computing a matrix inverse. This operation slightly restricts the scalability of OFT. Although we mitigate this issue by adopting a block diagonal parametrization, finding a faster and differentiable way to perform the matrix inverse remains a challenge. Second, OFT shows promising potential in compositionality: the orthogonal matrices obtained from multiple OFT finetuning tasks can be multiplied to yield another orthogonal matrix. Investigating whether this set of orthogonal matrices can retain the knowledge from all downstream tasks is an interesting avenue for future work. Lastly, the parameter efficiency of OFT heavily relies on the block diagonal structure, which introduces certain biases and limits model flexibility. Developing methods to improve parameter efficiency in a more effective and less biased manner remains a significant open problem.

\end{document}